\documentclass[11pt]{article}
\usepackage[margin=1in]{geometry}
\usepackage{amsmath,amssymb,amsthm,mathtools}
\usepackage{graphicx}
\usepackage{microtype}
\usepackage[hidelinks]{hyperref}

\newtheorem{theorem}{Theorem}
\newtheorem{lemma}[theorem]{Lemma}
\newtheorem{remark}[theorem]{Remark}
\newcommand{\N}{\mathbb N}
\newcommand{\C}{\mathbb C}
\newcommand{\cW}{\mathcal W}
\newcommand{\ip}[2]{\langle #1,#2\rangle}
\newcommand{\ran}{\operatorname{ran}}

\title{A Complete Minimal Hilbert System with a Doubly Universal Element}
\author{Candidate full solution packet\\\small arXiv:2306.11156, Questions 1 and 2}
\date{11 August 2026}

\begin{document}
\maketitle

\noindent\textbf{Status.} Candidate full affirmative answer to both questions,
pending expert review.  One and the same element for one and the same complete
minimal system is universal both by its ordinary partial sums and by
rearrangements.  Both the source and target spaces are separable
infinite-dimensional complex Hilbert spaces, and the target carries its usual
norm topology.

\section{Source questions}

Avetisyan, Grigoryan, and Ruzhansky ask the following on printed page 8 of
\cite{AGR}:

\begin{center}
\includegraphics[width=0.98\textwidth]{figures/open_problem_crop.png}
\end{center}

In their terminology, a Banach space $\mathcal X$ is continuously and densely
embedded in a Hausdorff sequential convergence space $\Omega$, and
$(\varphi_n)$ is a minimal system in $\mathcal X$ with biorthogonal
functionals $(c_n)$.  An element $U\in\mathcal X$ is universal in the usual
sense if
\[
  \left\{\sum_{n=1}^N c_n(U)\varphi_n:N\in\N\right\}
\]
has, for each $f\in\Omega$, a subsequence converging to $f$.  It is universal
in the sense of rearrangements if for every $f\in\Omega$ there is a
permutation $\sigma$ of $\N$ such that
\[
  \sum_{n=1}^\infty c_{\sigma(n)}(U)\varphi_{\sigma(n)}=f
  \quad\text{in }\Omega.
\]

The source calls a triple non-trivial when the topology of $\Omega$ is not too
weak.  The construction below takes $\Omega=\ell_2$ with its ordinary Hilbert
norm.  More strongly, $\mathcal X$ is itself a Hilbert space and
$(\varphi_n)$ is complete as well as minimal.

\section{Proof intuition}

There are three moving parts.  First, put in each coordinate of $\ell_2$ a
scaled positive and negative harmonic series, together with imaginary copies.
The scalar crossing algorithm can rearrange the four subseries in coordinate
$j$ to any prescribed complex number.  Their $2^{-j}$ scaling lets all
coordinate rearrangements be interleaved without losing $\ell_2$ convergence.
The same labeled terms can then be put in a fixed order whose ordinary partial
sums are universal.

Second, lift these target-space terms $w_n$ through the dense diagonal map
$B e_j=2^{-j}e_j$, writing $w_n=Bv_n$.  In a second Hilbert space set
\[
  \varphi_n=(e_n,a_n^{-1}v_n),\qquad a_n=2^{-n}.
\]
The first coordinate makes the coefficient functional simply
$c_n(x,y)=x_n$, so $U=(a,0)$ has the desired Fourier terms.  Completeness is
the key audit: an annihilator of all $\varphi_n$ would make
$\sum_n|\ip{v_n}{z}|$ finite, while the harmonic terms in any nonzero
coordinate of $z$ force it to diverge.

Finally, embed the two-coordinate Hilbert space into one copy of $\ell_2$ by
$I(x,y)=Ax+By$, where $A(H)$ is chosen algebraically disjoint from $B(H)$.
The $A$-part of the Fourier series converges unconditionally, while the
$B$-part is exactly the universal series $(w_n)$.  It follows at once that
ordinary partial sums and suitable rearrangements have the required limits.

\section{A series universal in both senses}

Throughout, $H=\ell_2(\N;\C)$ with standard basis $(e_j)$.  All members of
the families below are treated as \emph{labeled}; equal vectors, if any, are
still distinct terms.

\begin{lemma}[A rearrangement-universal Hilbert series]\label{lem:rearr}
For $j,m\geq1$, put $\delta_j=2^{-j}$ and consider the countable labeled
family
\[
 \cW=\left\{
  \frac{\delta_j}{m}e_j,-\frac{\delta_j}{m}e_j,
  \frac{i\delta_j}{m}e_j,-\frac{i\delta_j}{m}e_j:
  j,m\in\N
 \right\}.
\]
For every $x\in H$ there is an enumeration of all terms of $\cW$ whose series
converges in $H$ to $x$.
\end{lemma}

\begin{proof}
Fix $x=(x_j)\in H$.  In coordinate $j$, order the positive and negative
real terms $\pm\delta_j/m$ by the elementary crossing algorithm: take unused
positive terms until the partial sum first exceeds $\operatorname{Re}x_j$,
then unused negative terms until it first falls below that number, and repeat
(with the roles reversed if appropriate at the start).  Since both one-sided
harmonic sums diverge and their terms tend to zero, every term is eventually
used and the scalar partial sums converge to $\operatorname{Re}x_j$.  At all
times their absolute values are at most
$|\operatorname{Re}x_j|+\delta_j$.  Do the same with the imaginary terms and
interleave the two resulting scalar orders.  We obtain a local order
$(u_{j,r})_{r\geq1}$ of all four families in coordinate $j$ such that, for
\[
  s_{j,R}:=\sum_{r=1}^R u_{j,r},
\]
we have
\[
  s_{j,R}\longrightarrow x_j,
  \qquad
  |s_{j,R}|\leq C(|x_j|+\delta_j)                         \tag{1}
\]
with an absolute constant $C$.

We now interleave the local orders.  At stage $k$, for each $1\leq j\leq k$
in turn, append enough as-yet-unused terms of its local order so that at least
the first $k$ local terms have been used and so that all subsequent local
partial sums are within $2^{-k}\delta_j$ of $x_j$.  This is possible by local
convergence.  Every labeled term is eventually used.

It remains to check convergence between, not merely at, stage endpoints.
Given $\varepsilon>0$, choose $J$ so large that the $\ell_2$ tail of
$(|x_j|+\delta_j)$ beyond $J$ is smaller than
$\varepsilon/(4C+4)$.  Once the stage is sufficiently large, all local
partial sums in the first $J$ coordinates remain as close to $x_j$ as desired.
For $j>J$, whether a coordinate has not yet been started or is currently being
advanced, (1) bounds its error by
$(C+1)(|x_j|+\delta_j)$.  The first $J$ errors and the uniform tail bound are
therefore both small.  Thus every sufficiently late global partial sum is
within $\varepsilon$ of $x$ in $H$.
\end{proof}

\begin{lemma}[One order is also ordinarily universal]\label{lem:order}
The terms of $\cW$ admit an enumeration $(w_n)_{n\geq1}$ with both properties:
\begin{enumerate}
  \item the ordinary partial sums $\sum_{n=1}^Nw_n$ form a universal sequence
        in $H$;
  \item the series $\sum_nw_n$ is universal in $H$ under rearrangements.
\end{enumerate}
\end{lemma}

\begin{proof}
Deleting finitely many labeled terms preserves the conclusion of
Lemma~\ref{lem:rearr}.  Indeed, if the deleted family is $F$ and the desired
sum of the remaining family is $y$, rearrange the full family to
$y+\sum_{w\in F}w$ and delete the finitely many occurrences from that
convergent series.

Fix an initial numbering of all labels and choose a sequence $(d_k)$ in $H$
such that every point of $H$ is a subsequential limit; for example, repeat a
countable dense set in successive finite triangular blocks.  Suppose a finite
ordered set of labels has been used and its sum is $s$.  At stage $k$, first
append the least unused label, and update $s$.  The remaining cofinite family
is rearrangement universal, so it has an enumeration whose series sums to
$d_k-s$.  Append a finite initial segment of that enumeration which brings
the total sum within $2^{-k}$ of $d_k$.

The forced least-label step ensures that every label is eventually used, so
the infinite concatenation is an enumeration $(w_n)$ of $\cW$.  Its stage-end
partial sums approach the universal target sequence $(d_k)$ with errors
tending to zero; hence the ordinary partial sums are universal.  Property 2
depends only on the labeled family, not on its initial order, and follows from
Lemma~\ref{lem:rearr}.
\end{proof}

\section{A non-trivial dense Hilbert embedding}

\begin{lemma}\label{lem:embed}
There are bounded operators $A,B:H\to H$ such that $A$ is an isometry, $B$ is
injective with dense range, and
\[
  A(H)\cap B(H)=\{0\}.
\]
Consequently
\[
  I:H\oplus_2H\longrightarrow H,\qquad I(x,y)=Ax+By,
\]
is a bounded injective dense embedding whose inverse on its range is not
continuous.
\end{lemma}

\begin{proof}
Let $Be_k=2^{-k}e_k$.  This operator is bounded and injective, and its range
contains all finitely supported sequences, hence is dense.

Partition $\N$ into infinite disjoint sets $(S_j)_{j\geq1}$.  Choose
$\gamma_j>0$ so that
\[
  q_j:=\gamma_j\sum_{k\in S_j}2^{-k}e_k
\]
has norm one, and define $Ae_j=q_j$.  The vectors $(q_j)$ have disjoint
supports and are orthonormal, so $A$ is an isometry.

If $Ax=By$, then for every $k\in S_j$,
\[
  y_k=2^k(Ax)_k=\gamma_jx_j.
\]
Because $S_j$ is infinite, $y\in\ell_2$ forces $x_j=0$.  Thus $x=0$ and then
$y=0$, proving the disjointness and injectivity of $I$.  The range of $I$ is
dense because it contains $B(H)$.  Finally,
\[
  \|(0,e_k)\|_{H\oplus_2H}=1,
  \qquad
  \|I(0,e_k)\|_H=2^{-k}\longrightarrow0,
\]
so the induced target topology is strictly weaker than the source norm
topology.
\end{proof}

\section{The universal element}

\begin{theorem}\label{thm:main}
There exist separable infinite-dimensional complex Hilbert spaces
$\mathcal X$ and $\Omega$, a bounded injective dense embedding
$I:\mathcal X\to\Omega$, a complete minimal system
$(\varphi_n)_{n\geq1}$ in $\mathcal X$ with its unique biorthogonal system
$(c_n)_{n\geq1}$, and an element $U\in\mathcal X$ such that its Fourier
series
\[
  \sum_{n=1}^\infty c_n(U)\varphi_n
\]
is universal in $\Omega$ both in the usual sense and in the sense of
rearrangements.  The topology of $\Omega$ is its ordinary Hilbert norm and is
strictly weaker on the embedded copy of $\mathcal X$ than the norm topology
of $\mathcal X$.
\end{theorem}

\begin{proof}
Take $\Omega=H$, $\mathcal X=H\oplus_2H$, and the embedding $I$ from
Lemma~\ref{lem:embed}.  Fix the enumeration $(w_n)$ from
Lemma~\ref{lem:order}.  Every $w_n$ is supported in one coordinate and hence
belongs to $B(H)$; define
\[
  v_n:=B^{-1}w_n.
\]
As a labeled family, $(v_n)$ consists exactly of
\[
  \left\{
  \frac1m e_j,-\frac1m e_j,
  \frac{i}{m}e_j,-\frac{i}{m}e_j:j,m\in\N
  \right\}.                                                   \tag{2}
\]

Let $a_n=2^{-n}$, so $a=(a_n)\in H$, and put
\[
  \varphi_n=(e_n,a_n^{-1}v_n)\in\mathcal X,
  \qquad
  c_n(x,y)=x_n,
  \qquad
  U=(a,0).
\]
The $c_n$ are continuous and
$c_n(\varphi_m)=\delta_{nm}$, so $(\varphi_n)$ is minimal and $(c_n)$ is
biorthogonal to it.

We verify the stronger completeness property.  Suppose $(p,z)\in H\oplus_2H$
is orthogonal to every $\varphi_n$.  With the inner product taken linear in
the first variable,
\[
  0=\ip{\varphi_n}{(p,z)}
   =\ip{e_n}{p}+a_n^{-1}\ip{v_n}{z}.
\]
It follows from Cauchy--Schwarz that
\[
  \sum_{n=1}^\infty|\ip{v_n}{z}|
  =\sum_{n=1}^\infty a_n|\ip{e_n}{p}|
  \leq \|a\|_2\|p\|_2<\infty.                               \tag{3}
\]
If $z\ne0$, select $j$ with $\ip{e_j}{z}\ne0$.  The positive real
subfamily in (2) then gives
\[
  \sum_{m=1}^\infty
    \left|\ip{m^{-1}e_j}{z}\right|=\infty,
\]
contradicting (3).  Hence $z=0$, and the orthogonality equations give
$p=0$.  Thus the closed span of $(\varphi_n)$ is all of $\mathcal X$.
In particular the biorthogonal system is unique.

Since $c_n(U)=a_n$, the $N$th Fourier partial sum has image
\begin{align*}
 I\left(\sum_{n=1}^N c_n(U)\varphi_n\right)
 &=I\left(\sum_{n=1}^Na_ne_n,\sum_{n=1}^Nv_n\right)\\
 &=A a^{(N)}+\sum_{n=1}^Nw_n,                                \tag{4}
\end{align*}
where $a^{(N)}=(a_1,\ldots,a_N,0,\ldots)$.  Given $f\in H$, ordinary
universality in Lemma~\ref{lem:order} supplies a subsequence $N_k$ such that
\[
  \sum_{n=1}^{N_k}w_n\longrightarrow f-Aa.
\]
As $Aa^{(N_k)}\to Aa$, (4) converges to $f$.  This is universality in the
usual sense.

For rearrangement universality, fix $f\in H$ and choose a permutation
$\sigma$ such that
\[
  \sum_{n=1}^\infty w_{\sigma(n)}=f-Aa
\]
by Lemma~\ref{lem:order}.  The orthogonal series
$\sum_na_ne_n$ converges unconditionally to $a$, and hence
\[
 \sum_{n=1}^N I\bigl(c_{\sigma(n)}(U)\varphi_{\sigma(n)}\bigr)
 =A\sum_{n=1}^Na_{\sigma(n)}e_{\sigma(n)}
  +\sum_{n=1}^Nw_{\sigma(n)}
 \longrightarrow Aa+(f-Aa)=f.
\]
This is precisely universality in the sense of rearrangements.
Lemma~\ref{lem:embed} supplies all remaining embedding and non-triviality
claims.
\end{proof}

\section{Scope, verification, and novelty}

The theorem answers both existential questions exactly as stated, and it does
so with a complete minimal system rather than a merely minimal one.  It does
\emph{not} resolve the more concrete problem suggested by the authors, namely
whether one may take $\mathcal X=L^1(\mathcal M)$ and
$\Omega=L^p(\mathcal M)$ for $0\leq p<1$, nor does it produce a classical
trigonometric or Walsh system.  The construction is deliberately tailored.

The proof was audited at the following pressure points: the countable
interleaving converges at every partial sum rather than only at block
endpoints; a cofinite deletion really preserves rearrangement universality;
$A(H)\cap B(H)=\{0\}$; the annihilator argument gives completeness; and the
same permutation is compatible with the unconditionally convergent
first-coordinate series.  The accompanying \texttt{verification.md} records
these checks.  No numerical step is used as proof.

A bounded novelty search was performed on 11 August 2026.  The run registry,
solution, attempt, and proof-gap indexes were searched for arXiv:2306.11156,
the exact title, and the core phrases ``universal element'', ``minimal
system'', and ``universal by rearrangements''.  The locally parsed arXiv
corpus and a current arXiv-domain web search were searched for the exact
question and close phrases.  These searches recovered the source paper and
background work on universal \emph{series}, but no later paper answering the
element questions.  This is not an exhaustive bibliographic review, so
novelty remains subject to specialist confirmation.

\medskip
\noindent\textbf{Human-review recommendation.} Review as a candidate full
solution.  The highest-value checks are Lemma~\ref{lem:rearr}'s uniform tail
control, Lemma~\ref{lem:order}'s cofinite scheduling, and the intended scope of
the source phrase ``non-trivial triple.''

\begin{thebibliography}{9}
\bibitem{AGR}
Z. Avetisyan, M. Grigoryan, and M. Ruzhansky,
\emph{Elements (functions) that are universal with respect to a minimal
system}, arXiv:2306.11156 (2023), Questions 1 and 2 on printed page 8.
\end{thebibliography}

\end{document}
